\clearpage

\section{Supersolidity: Definition and
Conditions}\label{sec-nbcp-supersolidity}

\subsection{Density order and transverse phase
coherence}\label{density-order-and-transverse-phase-coherence}

For \(S=1/2\), the mapping \(n_i=\hat S_i^z+1/2\) and
\(b_i^\dagger=\hat S_i^+\) relates longitudinal spin order to
boson-density order and transverse spin coherence to boson phase
coherence. A three-sublattice density order parameter can be defined by

\begin{equation}\phantomsection\label{eq-nbcp-density-order}{
D_z=\frac13\left(\langle S_A^z\rangle+
e^{2\pi i/3}\langle S_B^z\rangle+
e^{-2\pi i/3}\langle S_C^z\rangle\right).
}\end{equation}

Nonzero \(D_z\) distinguishes a longitudinally modulated state from the
uniformly polarized state. A transverse ordered component is described
by \(\langle S_\alpha^+\rangle=\Psi_\alpha e^{i\phi}\) in a chosen
broken-symmetry state, with sublattice-dependent amplitudes
\(\Psi_\alpha\). The common phase is the relevant collective coordinate;
the sum of the three transverse moments can vanish even when these
amplitudes are nonzero.

At exact U(1) symmetry and zero temperature, coexistence of density
order and transverse long-range order supplies the usual spin-supersolid
criterion. At nonzero temperature in the strictly two-dimensional
phase-only regime, the transverse criterion is algebraic correlations
and a nonzero renormalized stiffness, while density order must persist.
A static Y/V-shaped arrangement, a large magnetocaloric response, or a
local energy minimum alone does not establish both properties.

\subsection{Spatial stiffness and continuum
normalization}\label{spatial-stiffness-and-continuum-normalization}

To separate local pinning from spatial rigidity, consider a continuum
energy with physical area element \(d^2r\),

\begin{equation}\phantomsection\label{eq-nbcp-continuum-energy}{
F[\phi]=\int d^2r\left[
\frac12\rho_{\mu\nu}\partial_\mu\phi\,\partial_\nu\phi
+\frac{v(\phi)}{a_{\rm spin}}\right],
\qquad a_{\rm spin}=\frac{\sqrt3a^2}{2}.
}\end{equation}

Here \(a\) is the nearest-neighbor distance, \(a_{\rm spin}\) is the
area per spin, \(\rho_{\mu\nu}\) has units of energy, and \(v\) is a
pinning energy per spin. After the other modes have been relaxed, the
coefficient of a small phase gradient defines the stiffness.
Uniform-rotation data determine \(v\) but contain no information about
that gradient coefficient. When the leading local dynamical reduction
has no mixed time-and-space derivative term, the small-wave-vector
dispersion satisfies

\begin{equation}\phantomsection\label{eq-nbcp-phase-dispersion}{
\varepsilon(\mathbf q)^2=
\frac{C_\phi+a_{\rm spin}\rho_{\mu\nu}q_\mu q_\nu}{\chi_z}.
}\end{equation}

The matching is tested for the zero-SOC classical Y state in
Section~\ref{sec-nbcp-y-stiffness}. With SOC, a mixed
frequency--momentum term can invalidate this reciprocal form even while
a static stiffness exists; Section~\ref{sec-nbcp-y-soc-conditions}
derives the replacement. The V stiffness and quantum and thermal
renormalizations remain undetermined.

\subsection{Clock RG and the algebraic window}\label{sec-nbcp-clock}

The clock description applies after short-distance modes have been
eliminated and the thermal coefficients have been matched. For a
constant positive stiffness tensor, use the effective stiffness
\(\rho_{\rm eff}=\sqrt{\det\rho}\) and dimensionless Gaussian
coefficient \(K=\rho_{\rm eff}/(k_{\rm B}T)\) at the relevant scale.
Here \(p\) is the clock harmonic, \(y_p\) its dimensionless pinning
amplitude, \(y_v\) the unit-vortex fugacity, and \(\eta\) the exponent
of transverse algebraic correlations. The field and fugacity
normalizations, shell integration and coupled flow are derived in
Section~\ref{sec-nbcp-clock-rg}.

At \(y_p=y_v=0\), simultaneous irrelevance requires \(2-p^2/(4\pi K)<0\)
and \(2-\pi K<0\). Thus

\begin{equation}\phantomsection\label{eq-nbcp-clock-window}{
\frac{2}{\pi}<K<\frac{p^2}{8\pi}.
}\end{equation}

A nonempty interval requires \(p>4\); \(p=4\) is marginal at the meeting
point and is not decided by these linear inequalities. For \(p=3\) the
interval is empty. For \(p=6\),

\begin{equation}\phantomsection\label{eq-nbcp-clock-sixfold-window}{
\frac{2}{\pi}<K<\frac{9}{2\pi},\qquad
\frac19<\eta<\frac14,\qquad
\frac{2\pi\rho_{\rm eff}}9<k_{\rm B}T<\frac{\pi\rho_{\rm eff}}2.
}\end{equation}

The last form is a condition using the stiffness at the relevant scale,
not a pair of transition temperatures obtained by inserting a classical
bare exchange. With nonzero initial fugacities, the trajectory must
reach the Gaussian basin. In a weak-clock scenario the low-temperature
side is pinned, the intermediate region has algebraic transverse
correlations, and the high-temperature side has proliferated vortices.
The limiting exponents \(1/9\) and \(1/4\) provide more direct
finite-size targets than an unrenormalized estimate of the transition
temperatures. Stable longitudinal density order must coexist with the
algebraic correlations for the intermediate state to meet the
spin-supersolid criterion used here.

\subsection{Microscopic tests of the clock
reduction}\label{sec-nbcp-clock-matching-tests}

The proposed clock description has a theoretical basis, but symmetry is
not sufficient to establish its applicability to NBCP. Sixfold pinning
can be irrelevant on a Gaussian fixed line, and numerical studies of the
standard ferromagnetic six-state clock model find two transitions; for
example, Chen et al.~(2017)~\cite{chen2017}
determine the lower and upper transitions with tensor renormalization.
These results establish properties of that effective model. They do not
show that a frustrated quantum magnet with additional density domains
and competing modes lies in the same basin of attraction. The symmetry
and universality argument in
Park et al., Sec. II~\cite{park2026} is
therefore a physically motivated effective-theory hypothesis, rather
than a completed microscopic finite-temperature derivation.

There is a controlled route in principle: start from a stable phase
theory, separate massive modes from the compact phase, and match all
relevant couplings before applying RG. In a weak-SOC regime continuously
connected to a U(1) algebraic phase, weak irrelevant perturbations
preserve that phase provided the density sector remains noncritical and
no additional relevant defect is introduced. This conditional argument
is stronger than a guess based only on the number of angular minima. Its
hypotheses must be checked for the chosen field and SOC strength. The
following tests specify the full program;
Section~\ref{sec-nbcp-y-full-zone-stability} supplies an initial
harmonic stability screen, not completion of the thermal matching.

\textbf{Separation of the other modes.} At each field and SOC point,
compute the constrained static Hessian and the dynamical spectrum over
the magnetic Brillouin zone, with mesh refinement near minima. Inspect
several points on the angular orbit and the selected minimum, allowing
the hard internal coordinates to relax. Track mode character with
eigenvectors or spectral projectors so that a band crossing is not
mistaken for the disappearance of a nonphase mode. An instability at
nonzero momentum, an extra soft mode or a nearly critical density
amplitude invalidates a single-phase reduction even if the uniform
five-dimensional block is positive.

Let \(\Delta_h\) be the lowest dynamical energy of a mode proposed for
elimination, and \(\xi_h\) its static correlation length. A local
low-energy elimination requires external scales
\(\varepsilon\ll\Delta_h\) and \(k\xi_h\ll1\). The inequality
\(k_BT\ll\Delta_h\) is a sufficient condition for freezing its thermal
occupation, not a necessary condition for integrating it out: thermally
occupied modes with finite short correlation lengths can instead
generate temperature-dependent coefficients. Raw eigenvalues of an
arbitrarily normalized static Hessian must not be compared with
temperature as though they were excitation energies. The output should
identify retained modes, eliminated modes, their scales and the momentum
range of the proposed phase theory. Local stability must also be
distinguished from a comparison with competing phases.

\textbf{Angular gradients and higher harmonics.} Compute the relaxed
tensor \(\rho_{\mu\nu}(\phi,T)\) across the full orbit from a
constrained small-amplitude spatial response, and compute the angular
free energy \(f_>(\phi,T)\) after eliminating the chosen fast modes. The
zero-temperature calculation already provides part of this starting
point. At finite temperature, a harmonic contribution contains
\(k_BT\sum\ln[1-\exp(-\varepsilon/k_BT)]\) as well as the properly
normalized vacuum energy. Low-energy phase modes retained for the
subsequent RG must not also be integrated out in this matching step. The
matching cutoff and energy normalization must be recorded.

Define unrestricted Fourier coefficients before imposing the expected
symmetry. For the uniform Y potential, write

\begin{equation}\phantomsection\label{eq-nbcp-matching-harmonics}{
f_>(\phi,T)=f_0(T)-\sum_{n=6,12,\ldots}
\lambda_n(T)\cos[n\phi-\delta_n(T)].
}\end{equation}

Forbidden harmonics provide symmetry and numerical diagnostics. Increase
the momentum mesh and angular resolution until the retained coefficients
and their uncertainties stabilize. Compare the contribution of the
higher harmonics to both the energy and its curvature, which weights the
\(n\)th amplitude by \(n^2\). A small \(\lambda_{12}/\lambda_6\) alone
is not a complete truncation test. Fourier-resolve the angle-dependent
gradient tensor as well, then compare the truncated and extended
effective theories under RG or finite-size statistical calculations.
Their predictions should be stable against adding harmonics and changing
the matching scale within an overlap region. The normal-ordered Gaussian
power counting given above supplies the weak-coupling expectation for
derivative terms; it does not replace this matching.

\textbf{Density walls and composite defects.} Keep the three
translation-related density domains as explicit boundary data. In a long
strip, impose different density domains on the two sides, vary the
transverse phase difference \(\Delta\phi\), and relax all spin
directions, including local canting and density amplitude. The excess
energy per wall length estimates a classical wall tension. At finite
temperature the required quantity is instead the excess constrained free
energy,

\begin{equation}\phantomsection\label{eq-nbcp-matching-wall-tension}{
\sigma(T,\Delta\phi)=\lim_{L_\parallel\to\infty}
\frac{F_{\rm wall}(T,\Delta\phi)-F_{\rm bulk}(T)}{L_\parallel}.
}\end{equation}

Boundary contributions must be subtracted consistently; a periodic
construction generally contains two walls and requires that factor to be
removed. Determine whether the density wall favors a nonzero phase
offset, whether it becomes soft, and whether its junctions bind phase
winding. A wall between density translation domains is not automatically
the same object as a wall between neighboring sixfold angular minima.

To measure defect binding and confinement, constrain a defect pair at
separation \(R\) and compare its relaxed energy or free energy with
isolated walls and vortex pairs. A useful diagnostic form, where a
phase-only description holds away from the cores, is

\begin{equation}\phantomsection\label{eq-nbcp-matching-defect-pair}{
F_{\rm pair}(R)-F_{\rm bulk}
\simeq 2F_{\rm core}
+2\pi\rho_{\rm eff}q^2\ln(R/a_0)+\sigma_{\rm string}R.
}\end{equation}

The core positions are held fixed here; translational entropy is
assigned later through the defect fugacity. The winding \(q\) and any
confining string must be derived from the full density-plus-phase order
parameter and its boundary conditions. Fractional vortices must not be
assumed to be allowed or deconfined. If a smaller winding is found to be
deconfined, its Gaussian fugacity eigenvalue is \(2-\pi Kq^2\), and the
integer-vortex clock criterion must be revised. A bound composite defect
can instead renormalize core weights without changing the smallest
asymptotically free winding.

\textbf{Matching and testing the RG trajectory.} At a common cutoff
\(a_0\), determine the smooth stiffness, angular potential and core free
energies in the same approximation and thermal ensemble. For a constant
positive tensor and unit vortices the initial variables are

\begin{equation}\phantomsection\label{eq-nbcp-matching-rg-inputs}{
K_0(T)=\frac{\sqrt{\det\rho_0(T)}}{k_BT},\qquad
y_{n,0}(T)=\frac{a_0^2\lambda_n(T)}{a_{\rm spin}k_BT},\qquad
y_{v,0}(T)=A_{\rm core}(T)e^{-F_{\rm core}(T)/(k_BT)}.
}\end{equation}

The core prefactor depends on the matching convention. The stiffness
here is the coefficient at the matching scale, not an already fully
renormalized thermodynamic stiffness to which the same vortex
corrections can be added again. Propagate these inputs and their
uncertainties through the coupled flow, including additional operators
if the earlier tests retain them. A weak-coupling RG is useful only
while its fugacities remain small. If they are large initially, a
nonperturbative statistical calculation of the matched model is needed.
An inferred algebraic regime should survive changes of matching cutoff
and truncation rather than appearing only for a particular initial
guess.

\phantomsection\label{sec-nbcp-matching-results}\textbf{First matching results.} Two
of these steps have been carried out for Y at 0.2 T,
\(J_{PD}=0.010\) meV and \(J_\Gamma=0\).

\emph{Matching cutoff.} Along the exact orbit the relaxed classical
susceptibility \(\chi\) is independent of \(\phi\) to \(10^{-15}\),
while the relaxed stiffness varies: \(\sqrt{\det\rho}=0.00238\) meV
with a 2.0\% \(\cos6\phi\) component. A phase theory with these
coefficients has the soft frequency
\(\varepsilon=k\,[\hat{\mathbf k}\cdot\rho(\phi)\cdot\hat{\mathbf k}/\chi]^{1/2}\).
Its sixfold harmonic of \(\langle\ln\varepsilon\rangle\) is 0.005034;
the LSWT soft branch extrapolated to \(k\to0\) gives 0.005022. Removing
the modes with \(|\mathbf k|<\Lambda\) from the harmonic thermal free
energy changes \(\lambda_6\); the phase theory reproduces that change
to 6--7\% at \(\Lambda=0.25\) for \(T=0.0025\)--\(0.02\) meV, and
misses it by 25--30\% at \(\Lambda=0.5\), outside its range. The
dependence on the matching cutoff is therefore not an ambiguity. Pinning
removed from the matching is generated again, at the same order, when
the phase theory integrates its own modes with the angle-dependent
stiffness. A matching with a cutoff must keep \(\rho(\phi)\); one with
a constant stiffness must keep all modes in \(f_>\). Both give the same
long-wavelength coefficient at \(O(T)\), which for Y is
\(\lambda_6(T)/\lambda_6(0)=0.90\), 0.67 and 0.62 at \(T=0.005\),
0.010 and 0.015 meV. The order-by-disorder entropy of the anisotropic
stiffness opposes the zero-point selection. A larger cutoff, such as
\(\Lambda=1\), removes modes the phase theory cannot represent, and its
change of sign is an artefact of that split. Not included are the
Debye--Waller suppression by the phase fluctuations themselves, which
is the RG eigenvalue and is of relative order \(T/S\) beyond this,
anharmonic terms and defects.

\emph{Vortex cores.} Relaxed classical vortices in hexagonal clusters
of up to 6\,937 spins, with two outer rings fixed to the winding,
give \(E_v(R)=E_{\rm core}+\kappa\ln R\). At \(J_{PD}=0\),
\(\kappa=0.01189\) meV against \(\pi\sqrt{\det\rho}=0.01201\) meV, and
\(E_{\rm core}=0.0071\)--0.0076 meV for both windings. At
\(J_{PD}=0.010\) meV, the winding \(m=-1\) gives \(\kappa=0.00719\)
against \(0.00763\) meV and \(E_{\rm core}=0.0036\)--0.0053 meV. The
ranges compare the fitted constant with the constant at the predicted
\(\kappa\). The vortex energy thus confirms independently that
\(J_{PD}=0.010\) meV lowers \(\sqrt{\det\rho}\) by about 37\%, which
lowers the bare upper window edge \(\pi\rho/2\) from 0.0060 to
0.0038 meV. With \(J_{PD}\neq0\) the two windings are inequivalent,
because the spin-only reflection that maps \(m\to-m\) at
\(J_{PD}=0\) is broken. The \(m=+1\) energy contains a term linear in
\(R\) whose sign follows the far-field angle. It is the boundary value
of a total-derivative gradient term, which vanishes for a uniform twist
but not for a co-rotating winding. In a neutral pair it acts only near
the cores, so its division between core and boundary is a convention;
only pair energies are free of it, and they are not yet available. Near
the window, \(E_{\rm core}/k_BT\approx1\)--2 for \(m=-1\), so
\(y_{v,0}\approx0.1\)--0.3. The vortex fugacity is not small, and a
weak-coupling flow from the bare stiffness overestimates the transition
temperature. This is consistent with the classical Monte Carlo
\(T_{\rm BKT}\approx0.0024\) meV of the roadmap calculation, which lies a
factor 1.6 below the reduced bare edge; for the lattice XY model the
corresponding factor is 1.76. That Monte Carlo stiffness estimator
omits the \(J_{PD}\) part of the exchange, so the comparison is
qualitative. These are classical cores; the quantum
\(S=1/2\) cores and the temperature dependence of \(E_{\rm core}\) are
not computed.

Finally test the thermal theory on increasing sizes. The target
observations are transverse correlations
\(G_\perp(r)\propto r^{-\eta}\), loss of sixfold locking in the
long-distance phase distribution, a finite appropriately defined phase
stiffness, and persistent longitudinal density order. For the simple
sixfold scenario the algebraic exponent lies between \(1/9\) and
\(1/4\), with substantial finite-size corrections near the boundaries.
The phase-locking diagnostic can be defined as
\(W_6(L)=\langle\cos[6\arg\Psi-\delta_6]\rangle\), where \(\Psi\) is the
transverse order parameter projected onto the chosen Y ordering pattern.
Its size dependence, together with correlations and density-order
observables, is more informative than observing six minima in a
zero-temperature energy plot.

The efficient order is spectrum and mode separation, then angular
matching, then real-space defects, followed by the thermal flow and size
scaling. A classical spin simulation tests a classical approximation; a
simulation of the matched phase theory tests the matching assumptions.
Neither alone constitutes a calculation of the full \(S=1/2\) quantum
magnet. Quantitative quantum conclusions require a controlled assessment
of the semiclassical error or an independent quantum thermal method with
spatial-size convergence. The four angular panels in
Fig.~\ref{fig-yv-orbit} complete only the zero-temperature potential
comparison, not these thermal tests.
